% !TEX root = ../main.tex
\subsection{Dataset}
\label{sec:dataset}

\textsc{SciFig-Eval} comprises 250 English-language scientific figures extracted from 187 arXiv papers published between 2023 and 2025, spanning primarily NLP, machine learning, and computational linguistics. We sample three chart types that dominate scientific reporting: bar charts (99), line plots (99), and pie charts (52). Each figure is paired with a structured expert annotation describing axes, data series, labels, color mappings, and trends. Two trained annotators with graduate-level NLP expertise produced the annotations, achieving 94\% agreement on verifiable checklist items (exact values, label names, axis ranges). Disagreements were resolved by a third adjudicator consulting the source PDF. We additionally record source page numbers and verified figure numbers, enabling in-paper evaluation where the model receives the full PDF page. We prioritize annotation depth over corpus scale; each of the 100 primary figures receives evaluation across 14 experimental conditions and 8~models, yielding over 11,000 individual evaluation instances.

\paragraph{Evaluation subsets.}
From the 250-figure pool, we construct a primary subset of 100 figures (40 bar, 40 line, 20 pie) via stratified random sampling (seed$=$42), and a nested ablation subset of 50 figures (20/20/10) for probe-designer independence studies (\S\ref{sec:ablation}).

\paragraph{Adversarial stimuli.}
We construct four categories of adversarial stimuli targeting distinct behavioral dimensions (defined in \S\ref{sec:ari}). \emph{Resistance probes} (250 figures, three types per figure) test whether models push back on false premises. \emph{Caption bias probes} (100 figures) pair each figure with a modified caption embedding 2--3 plausible but incorrect claims. \emph{Admittance blur} (50 figures) selectively obscures elements whose identity cannot be recovered from context. \emph{Inductance blur} (48 figures) obscures elements that remain inferable from surrounding visual cues. Blur candidates were identified through multi-model consensus and confirmed via human review on a purpose-built dashboard. Full construction details appear in Appendix~\ref{sec:dataset-details}.

% ── Table 1: Dataset statistics ──
\begin{table}[t]
\centering
\small
\begin{tabular}{@{}lr@{}}
\toprule
\textbf{Component} & \textbf{Count} \\
\midrule
\multicolumn{2}{@{}l}{\emph{Figure corpus}} \\
\quad Bar charts & 99 \\
\quad Line plots & 99 \\
\quad Pie charts & 52 \\
\quad Total figures & 250 \\
\midrule
\multicolumn{2}{@{}l}{\emph{Evaluation subsets}} \\
\quad Primary subset (40/40/20) & 100 \\
\quad Ablation subset (20/20/10) & 50 \\
\midrule
\multicolumn{2}{@{}l}{\emph{Adversarial stimuli}} \\
\quad Resistance probes & 250 \\
\quad Caption bias probes & 100 \\
\quad Admittance blur candidates & 50 \\
\quad Inductance blur candidates & 48 \\
\bottomrule
\end{tabular}
\caption{Dataset statistics for \textsc{SciFig-Eval}. The figure corpus spans three chart types drawn from arXiv papers. Evaluation subsets are stratified samples (seed$=$42). Adversarial stimuli target distinct behavioral dimensions of the A-R-I framework.}
\label{tab:dataset}
\end{table}
